\documentclass[UTF8,11pt]{ctexart}
\usepackage[a4paper,margin=24mm]{geometry}
\usepackage{amsmath,amssymb,amsthm,mathtools,mathrsfs}
\usepackage[all]{xy}
\usepackage{xurl,hyperref,enumitem}
\usepackage{tikz}
\usetikzlibrary{arrows.meta,cd,decorations.markings,calc,patterns}
\hypersetup{hidelinks,breaklinks=true}
\setlength{\emergencystretch}{3em}
\allowdisplaybreaks
\newtheorem*{theorem}{Theorem}
\newtheorem*{thm}{Theorem}
\newtheorem*{lemma}{Lemma}
\newtheorem*{proposition}{Proposition}
\newtheorem*{prop}{Proposition}
\newtheorem*{corollary}{Corollary}
\newtheorem*{cor}{Corollary}
\newtheorem*{claim}{Claim}
\theoremstyle{definition}
\newtheorem*{definition}{Definition}
\newtheorem*{defn}{Definition}
\newtheorem*{remark}{Remark}
\newtheorem*{example}{Example}
\newcommand{\R}{\mathbb R}
\newcommand{\C}{\mathbb C}
\newcommand{\Z}{\mathbb Z}
\newcommand{\Q}{\mathbb Q}
\newcommand{\N}{\mathbb N}
\newcommand{\F}{\mathbb F}
\newcommand{\D}{\mathbb D}
\newcommand{\bB}{\mathbb B}
\newcommand{\bC}{\mathbb C}
\newcommand{\bR}{\mathbb R}
\newcommand{\bZ}{\mathbb Z}
\newcommand{\bN}{\mathbb N}
\newcommand{\bQ}{\mathbb Q}
\newcommand{\bD}{\mathbb D}
\newcommand{\bF}{\mathbb F}
\newcommand{\CP}{\mathbb{CP}}
\newcommand{\RP}{\mathbb{RP}}
\newcommand{\sO}{\mathscr O}
\newcommand{\sI}{\mathscr I}
\newcommand{\sF}{\mathscr F}
\newcommand{\sG}{\mathscr G}
\newcommand{\sH}{\mathscr H}
\newcommand{\sR}{\mathscr R}
\newcommand{\sM}{\mathscr M}
\newcommand{\sV}{\mathscr V}
\newcommand{\sU}{\mathscr U}
\DeclareMathOperator{\Hom}{Hom}
\DeclareMathOperator{\Ext}{Ext}
\DeclareMathOperator{\Tor}{Tor}
\DeclareMathOperator{\Spec}{Spec}
\DeclareMathOperator{\Proj}{Proj}
\DeclareMathOperator{\supp}{supp}
\DeclareMathOperator{\im}{im}
\DeclareMathOperator{\id}{id}
\DeclareMathOperator{\Tr}{Tr}
\DeclareMathOperator{\tr}{tr}
\DeclareMathOperator{\rank}{rank}
\DeclareMathOperator{\ord}{ord}
\DeclareMathOperator{\dist}{dist}
\DeclareMathOperator{\End}{End}
\DeclareMathOperator{\Aut}{Aut}
\DeclareMathOperator{\Gal}{Gal}
\DeclareMathOperator{\vol}{vol}
\DeclareMathOperator{\Cl}{Cl}
\DeclareMathOperator{\coker}{coker}
\DeclareMathOperator{\codim}{codim}
\DeclareMathOperator{\divergence}{div}
\DeclareMathOperator{\Ric}{Ric}
\DeclareMathOperator{\grad}{grad}
\DeclareMathOperator{\Hess}{Hess}
\newcommand{\abs}[1]{\left\lvert#1\right\rvert}
\newcommand{\sabs}[1]{\lvert#1\rvert}
\newcommand{\babs}[1]{\bigl\lvert#1\bigr\rvert}
\newcommand{\Babs}[1]{\Bigl\lvert#1\Bigr\rvert}
\newcommand{\bbabs}[1]{\biggl\lvert#1\biggr\rvert}
\newcommand{\BBabs}[1]{\Biggl\lvert#1\Biggr\rvert}
\newcommand{\norm}[1]{\left\lVert#1\right\rVert}
\newcommand{\snorm}[1]{\lVert#1\rVert}
\newcommand{\bnorm}[1]{\bigl\lVert#1\bigr\rVert}
\newcommand{\Bnorm}[1]{\Bigl\lVert#1\Bigr\rVert}
\newcommand{\bbnorm}[1]{\biggl\lVert#1\biggr\rVert}
\newcommand{\BBnorm}[1]{\Biggl\lVert#1\Biggr\rVert}
\newcommand{\widebar}[1]{\overline{#1}}
\newcommand{\nicefrac}[2]{\frac{#1}{#2}}
\newcommand{\linnprod}[2]{\langle#1,#2\rangle}
\newcommand{\myindex}[1]{#1}
\newcommand{\glsadd}[1]{}
\newcommand{\avoidbreak}{}
\newcommand{\sourceref}[1]{\textnormal{[原文：\nolinkurl{#1}]}}
\newcommand{\thmref}[1]{\sourceref{#1}}
\newcommand{\lemmaref}[1]{\sourceref{#1}}
\newcommand{\propref}[1]{\sourceref{#1}}
\newcommand{\corref}[1]{\sourceref{#1}}
\newcommand{\defnref}[1]{\sourceref{#1}}
\newcommand{\exampleref}[1]{\sourceref{#1}}
\newcommand{\exerciseref}[1]{\sourceref{#1}}
\newcommand{\figureref}[1]{\sourceref{#1}}
\newcommand{\sectionref}[1]{\sourceref{#1}}
\newcommand{\appendixref}[1]{\sourceref{#1}}
\newcommand{\stacksref}[2]{\href{https://stacks.math.columbia.edu/tag/#1}{#2 [#1]}}


\newcommand{\E}{\mathbb E}
\newcommand{\Prob}{\mathbb P}
\newcommand{\Reff}{\mathcal R_{\mathrm{eff}}}
\newcommand{\eps}{\varepsilon}
\DeclareMathOperator{\diam}{diam}
\DeclareMathOperator{\Var}{Var}
\DeclareMathOperator{\Crit}{Crit}
\newcommand{\dd}{\,\mathrm d}


\begin{document}
\section*{079\quad 一阶逻辑的Craig插值与共享语言模型拼合}
\subsection*{来源与整理范围}
Open Logic Project, \emph{The Open Logic Text}, model-theory/interpolation：\texttt{separation.tex} (blob 16e6eaefbb367ae4dc4d03048fb6abdef92f0306) 与 \texttt{interpolation-proof.tex} (blob 021b8648fb8616eab1ca5af46490af7634976915)。收录两个分离引理及主定理的全部模型构造路线。
\par\url{https://github.com/OpenLogicProject/OpenLogic/blob/master/content/model-theory/interpolation/interpolation-proof.tex}
\par\textbf{恢复方式（整理者说明）：}依据人类原文重新整理以补齐缺失题号；并非旧题证文件的逐字节恢复；2026-09-28。
\subsection*{作者命题与人类证明}

\paragraph{符号展开说明（整理者）。}
以下以 $\mathcal L_i$ 表示语言，$\models$ 表示语义后承，$S[c/x]$ 表示代入。
Open Logic 的词汇宏和公式宏展开为常规记法。取通常的非空论域语义。

\begin{definition}[Separation]
A sentence $C$ separates sets of sentences $\Gamma$ and $\Delta$ if
$\Gamma\models C$ and $\Delta\models\neg C$. If no such sentence
in the specified common language exists, they are inseparable.
\end{definition}
\begin{lemma}[Adding new constants; separation, Lemma sep1]
Suppose $\mathcal L_0$ contains the nonlogical symbols common to $\Gamma$
and $\Delta$, and $\mathcal L'_0$ is obtained by adjoining infinitely many
new constants $c_n$. If $\Gamma$ and $\Delta$ are inseparable in
$\mathcal L_0$, they are inseparable in $\mathcal L'_0$.
\end{lemma}
\begin{proof}
Suppose instead they are separated in $\mathcal L'_0$.
Then $\Gamma\models C[c/x]$ and $\Delta\models\neg C[c/x]$ for some
$\mathcal L_0$-formula $C$, with $c$ a fresh constant; the case of several
new constants is similar. By compactness there are finite
$\Gamma_0\subseteq\Gamma$ and $\Delta_0\subseteq\Delta$ with the same
consequences. Let $G$ and $H$ be their respective conjunctions. Then
$G\models C[c/x]$ and $H\models\neg C[c/x]$.
Generalization gives $G\models\forall x\,C$. Contraposition of the second
consequence gives $C[c/x]\models\neg H$, and hence
$\forall x\,C\models\neg H$. Thus $H\models\neg\forall x\,C$.
By monotonicity, $\forall x\,C$ separates $\Gamma$ and $\Delta$
in $\mathcal L_0$, a contradiction.
\end{proof}
\begin{lemma}[Adding existential witnesses; separation, Lemma sep2]
Suppose $\Gamma\cup\{\exists x\,S\}$ and $\Delta$ are inseparable, and
$c$ is a new constant not in $\Gamma,\Delta,S$.
Then $\Gamma\cup\{\exists x\,S,S[c/x]\}$ and $\Delta$ are inseparable.
\end{lemma}
\begin{proof}
Suppose $C$ separates the latter pair. There are two cases.
If $c$ does not occur in $C$, the theory
$\Gamma\cup\{\exists x\,S,\neg C\}$ is satisfiable; otherwise $C$
already separates the original pair. It remains satisfiable after
adding $S[c/x]$, contrary to the claimed separation.
If $c$ occurs in $C$, write the separator as $C[c/x]$. We have
\[
\Gamma,\exists x\,S,S[c/x]\models C[c/x].
\]
Deduction and generalization give
$\Gamma,\exists x\,S\models\forall x(S\to C)$, and therefore
$\Gamma,\exists x\,S\models\exists x\,C$.
On the other hand $\Delta\models\neg C[c/x]$ gives
$\Delta\models\neg\exists x\,C$ by generalization.
The original pair is thus separable, again a contradiction.
\end{proof}

\begin{theorem}[Craig's Interpolation Theorem]
If $\models A\to B$, there is a sentence $C$ such that
$\models A\to C$ and $\models C\to B$, and every constant, function,
and predicate other than equality occurring in $C$ occurs in both $A$
and $B$.
\end{theorem}
\begin{proof}
Let $\mathcal L_1$ be the language of $A$, $\mathcal L_2$ that of $B$,
and $\mathcal L_0=\mathcal L_1\cap\mathcal L_2$. Adjoin the same infinitely
many new constants $c_0,c_1,\ldots$ to all three languages, obtaining
$\mathcal L'_i$.
If $A$ is unsatisfiable, $\exists x(x\ne x)$ is an interpolant.
If $\neg B$ is unsatisfiable, $\exists x(x=x)$ is an interpolant.
We may assume both $A$ and $\neg B$ are satisfiable.

We prove the contrapositive. Assume there is no interpolant, so
$\{A\}$ and $\{\neg B\}$ are inseparable in $\mathcal L_0$.
Our goal is to extend this pair to a maximally inseparable pair
$(\Gamma^*,\Delta^*)$.
Enumerate the sentences of $\mathcal L'_1$ as $A_0,A_1,\ldots$ and those
of $\mathcal L'_2$ as $B_0,B_1,\ldots$.
Set $\Gamma_0=\{A\}$ and $\Delta_0=\{\neg B\}$.
Define increasing sequences as follows.
\begin{enumerate}
\item If $\Gamma_n\cup\{A_n\}$ and $\Delta_n$ are inseparable in
$\mathcal L'_0$, put $A_n$ in $\Gamma_{n+1}$; if $A_n=\exists x\,S$,
also put $S[c/x]$ there, choosing $c$ not occurring in
$\Gamma_n,\Delta_n,A_n,B_n$. Otherwise leave $\Gamma_n$ unchanged.
\item If $\Gamma_{n+1}$ and $\Delta_n\cup\{B_n\}$ are inseparable,
put $B_n$ in $\Delta_{n+1}$; if $B_n=\exists x\,S$, also put
$S[c/x]$ there, choosing $c$ fresh for
$\Gamma_{n+1},\Delta_n,A_n,B_n$. Otherwise leave $\Delta_n$ unchanged.
\end{enumerate}
Finally let $\Gamma^*=\bigcup_n\Gamma_n$ and
$\Delta^*=\bigcup_n\Delta_n$.
By simultaneous induction,
\begin{enumerate}[label=(\alph*)]
\item\label{craig-a} $\Gamma_n$ and $\Delta_n$ are inseparable;
\item\label{craig-b} $\Gamma_{n+1}$ and $\Delta_n$ are inseparable.
\end{enumerate}
The basis for \ref{craig-a} follows from the first lemma. For
\ref{craig-b}, if nothing is added it is just \ref{craig-a}; if only
$A_n$ is added it holds by construction; if a witness is added as well,
the second lemma applies to $\Gamma_n\cup\{\exists x\,S\}$ and $\Delta_n$.
For the next instance of \ref{craig-a}, either $\Delta_{n+1}=\Delta_n$,
in which case use \ref{craig-b}, or $B_n$ is added, with inseparability
ensured by construction and by the second lemma when necessary.
The next instance of \ref{craig-b} follows in exactly the same three
cases for $A_{n+1}$. This completes the simultaneous induction.

It follows by compactness that $\Gamma^*$ and $\Delta^*$ are inseparable:
a separator would separate some finite stage. In particular both are
consistent, since an inconsistent member is separated from the other
by a contradictory or a valid sentence.

We show $\Gamma^*$ is maximally consistent in $\mathcal L'_1$, and
similarly for $\Delta^*$ in $\mathcal L'_2$.
Suppose both $A_n$ and $\neg A_n$ are absent from $\Gamma^*$.
Failure to add $A_n$ supplies a common-language sentence $C$ with
\[
\Gamma^*\models A_n\to C,\qquad \Delta^*\models\neg C.
\]
Failure to add $\neg A_n$ at its enumerated stage likewise supplies $C'$ with
\[
\Gamma^*\models\neg A_n\to C',\qquad \Delta^*\models\neg C'.
\]
Propositional logic gives $\Gamma^*\models C\lor C'$ and
$\Delta^*\models\neg(C\lor C')$, a contradiction.
The argument for $\Delta^*$ is the same.

Next, $\Gamma^*\cap\Delta^*$ is maximally consistent in $\mathcal L'_0$.
It is consistent as an intersection of consistent sets.
For any common-language sentence $S$, each of $\Gamma^*,\Delta^*$ contains
$S$ or its negation. Opposite choices would make $S$ or $\neg S$ a
separator. Thus either both contain $S$ or both contain $\neg S$.

By the witness construction, just as in the Henkin completeness proof,
$\Gamma^*$ has a model $M'_1$ whose domain consists of the interpretations
$c^{M'_1}$ of the constants; likewise $\Delta^*$ has a model $M'_2$ of
this form. Drop the symbols outside $\mathcal L'_0$ to obtain reducts
$M_1,M_2$. The map
\[
h:M_1\longrightarrow M_2,\qquad h(c^{M'_1})=c^{M'_2}
\]
is an isomorphism in $\mathcal L'_0$. Indeed each sentence in this
language belongs to both maximally consistent theories or to neither.
Equalities of constants therefore show well-definedness and injectivity;
all elements are named, so the map is onto. For predicates, for example,
\[
c^{M'_1}\in P^{M'_1}\quad\Longleftrightarrow\quad
P(c)\in\Gamma^*\quad\Longleftrightarrow\quad
P(c)\in\Delta^*\quad\Longleftrightarrow\quad c^{M'_2}\in P^{M'_2}.
\]
The other isomorphism conditions are established similarly.

Construct a model $M$ for the union language on the domain of $M'_2$:
\begin{enumerate}
\item Predicates belonging only to $\mathcal L'_2$ retain their
interpretation in $M'_2$.
\item Predicates belonging only to $\mathcal L'_1$ are transported by
$h$: $P^M=h(P^{M'_1})$, with $h$ applied coordinatewise.
\item For a common predicate the two prescriptions agree, since
$P^{M'_2}=h(P^{M'_1})$.
\item Functions, including constants, are handled in the same way.
\end{enumerate}
Induction on formulas shows that $M$ agrees with $M'_1$ on the first
language and with $M'_2$ on the second. Hence
$M\models\Gamma^*\cup\Delta^*$, so $M\models A\land\neg B$.
Thus $\not\models A\to B$, proving the contrapositive.
\end{proof}

\subsection*{整理说明与先修范围（非作者正文）}
允许引用一阶紧致性、带见证的Henkin模型构造和满足关系的同构不变性。分离保持、两套见证扩张、公共语言上的同构和合并解释都是本题核心，均保留。公式宏展开，少量重复归纳语句合并，但不省略归纳分支。原文枚举处写$\mathcal L_i$而后文要求$\mathcal L'_i$，这里按其新增常量约定统一；分离引理的$\neg\delta$统一为同段$\neg H$；模型拼合第三项误写$h(P^{M'_2})$，按紧随的坐标条件写为$h(P^{M'_1})$。最后模型使用扩张后的联合语言再取原语言还原。上述为明列的记号/指标整理，不另补其他证明路线。难度在于同步保持不可分离性并构造可拼合的见证模型，不是直接应用完备性。
\subsection*{可修改的简短标注（非作者正文）}
$c$：一阶语义后承的公共词汇插值

$a$：不可分离理论对；Henkin见证；紧致性；公共约化同构；解释运输

\texttt{cross\_theory\_status = yes}。主证明将找不到公共语言公式的句法限制转为两套模型的共同约化，通过常量解释构成同构并运输另一语言的结构，从而构造原后承关系的反模型。

全部标注为非穷尽、可修改初稿。
\subsection*{来源许可与编辑归属}
Open Logic Project，\emph{The Open Logic Text}，CC BY 4.0；许可见仓库README：\url{https://github.com/OpenLogicProject/OpenLogic/blob/master/README.md}。正文来源宏展开为标准LaTeX；编辑题号、取材范围与标注另列。
ChatGPT 加入题号、中文来源说明、整理范围和初稿标注；编辑说明与人类证明分列。
\end{document}
